\newcommand*{\ollangid}{bn}
\newcommand*{\ollanguage}{bengali}
\usepackage{fontspec}
\setmainfont{Latin Modern Roman}
\setsansfont{Latin Modern Sans}
\setmonofont{Latin Modern Mono}
\newfontfamily\latinfont{Latin Modern Roman}
% Character classes follow the accepted Bengali OpenLogic edition.
\newfontfamily\bengalifont{NotoSerifBengali-Regular.ttf}[
  Path=../../../../fonts/,Script=Bengali,Language=Bengali,
  BoldFont=NotoSerifBengali-Bold.ttf,
  ItalicFont=NotoSerifBengali-Regular.ttf,BoldItalicFont=NotoSerifBengali-Bold.ttf,
  ItalicFeatures={FakeSlant=0.15},BoldItalicFeatures={FakeSlant=0.15},
  SmallCapsFont=NotoSerifBengali-Regular.ttf]
\usepackage[Bengali,DevanagariDanDa,BasicLatin,GeneralPunctuation]{ucharclasses}
\XeTeXcharclass"200C=\BengaliClass
\XeTeXcharclass"200D=\BengaliClass
\XeTeXgenerateactualtext=1
% A character run can end outside a TeX group; do not open groups in hooks.
\setTransitionsFor{Bengali}{\bengalifont}{\latinfont}
\setTransitionsFor{DevanagariDanDa}{\bengalifont}{\latinfont}
\errorcontextlines=40
\definecolor{OLPcolor}{RGB}{151,86,7}
\definecolor{OLPdkcolor}{RGB}{104,65,6}
\definecolor{OLPltcolor}{RGB}{248,203,137}
\usepackage[nonumberlist,toc,style=index]{glossaries}
\makeglossaries
